\chapter{Esercizio sulla diffrazione di raggi X e introduzione alla dinamica reticolare}



Il professore fornisce una lista di esercizi sulla diffrazione di raggi X tratti da esami e prove in itinere degli anni precedenti, disponibili con soluzioni sul suo sito web:

22 gennaio 2019 | 8 aprile 2019 | 18 giugno 2019 | 10 settembre 2019 | 19 novembre 2019 | 3 febbraio 2020 | 15 giugno 2020



\section{Prova del 19 novembre 2018}

\begin{quote}
Si consideri un reticolo di Bravais con \textbf{simmetria cubica} e formula chimica \textbf{AB}. Utilizzando il metodo di \textbf{Debye-Scherrer} con lunghezza d'onda $\lambda = 0{,}1542$ nm, si osservano i primi 10 picchi di diffrazione ai seguenti angoli $\varphi = 2\theta$ (in gradi):
\end{quote}

\begin{table}[h!]
\centering
\begin{tabular}{|c|c|}
\hline
Picco & $2\theta$ [deg] \\ \hline
1 & $36{,}95$ \\
2 & $42{,}91$ \\
3 & $62{,}30$ \\
4 & $74{,}64$ \\
5 & $78{,}64$ \\
6 & $94{,}06$ \\
7 & $105{,}75$ \\
8 & $109{,}78$ \\
9 & $127{,}29$ \\
10 & $143{,}77$ \\ \hline
\end{tabular}
\end{table}

\textbf{Domanda 1:} Dopo aver dimostrato che il reticolo dato è un FCC, etichettare ciascun picco con gli indici di Miller $(hkl)$ del cubico semplice e calcolare il parametro reticolare $a$ della cella convenzionale.

\textbf{Domanda 2:} Indicare quali dei picchi sopra elencati scomparirebbero se gli atomi A e B fossero identici e l'atomo B fosse spostato di un vettore $\mathbf{d}_2 = a\left(\frac{1}{4}, \frac{1}{4}, \frac{1}{4}\right)$ oppure $\mathbf{d}_2 = a\left(\frac{1}{2}, \frac{1}{2}, \frac{1}{2}\right)$.

\subsection{Determinazione del reticolo di Bravais}

Si ricorda che nel metodo di Debye-Scherrer il modulo del vettore del reticolo reciproco è legato all'angolo di scattering $\varphi = 2\theta$ dalla relazione:
\begin{equation}
|\mathbf{K}| = 2|\mathbf{k}|\sin\theta
\end{equation}
dove $|\mathbf{k}| = \frac{2\pi}{\lambda}$ è il modulo del vettore d'onda incidente e $\theta = \varphi/2$ è il semiangolo di scattering (angolo di Bragg).

Per determinare il tipo di reticolo non è necessario conoscere il valore assoluto di $|\mathbf{K}|$: è sufficiente calcolare i \textbf{rapporti} tra i moduli dei vettori del reticolo reciproco, poiché:
\begin{equation}
\frac{K_n}{K_1} = \frac{\sin\theta_n}{\sin\theta_1}
\end{equation}


Si calcola $\theta = (2\theta)/2$ per ciascun picco, poi $\sin\theta$ e infine il rapporto $K_n/K_1$:

\begin{table}[h!]
\centering
\begin{tabular}{|c|c|c|c|c|}
\hline
Picco & $\theta$ [deg] & $\sin\theta$ & $K/K_1$ & Valore teorico BCC \\ \hline
1 & $18{,}475$ & $0{,}3169$ & $1$ & $1$ \\
2 & $21{,}455$ & $0{,}3658$ & $1{,}154$ & $\frac{2}{\sqrt{3}} \approx 1{,}155$ \\
3 & $31{,}150$ & $0{,}5173$ & $1{,}632$ & $\frac{2\sqrt{2}}{\sqrt{3}} \approx 1{,}633$ \\
4 & $37{,}320$ & $0{,}6063$ & $1{,}913$ & $\sqrt{\frac{11}{3}} \approx 1{,}915$ \\ \hline
\end{tabular}
\end{table}



I rapporti calcolati corrispondono esattamente ai rapporti attesi per un \textbf{reticolo reciproco BCC}. In un BCC, i vettori del reticolo reciproco ordinati per lunghezza crescente sono:

\begin{enumerate}
    \item Il vettore dal vertice al \textbf{centro del corpo} del cubo (il più corto): lunghezza $\frac{\sqrt{3}}{2} \cdot \frac{4\pi}{a}$, normalizzata a $1$.
    \item Il \textbf{lato} del cubo: rapporto $\frac{2}{\sqrt{3}}$.
    \item La \textbf{diagonale di faccia}: rapporto $\frac{2\sqrt{2}}{\sqrt{3}}$.
    \item Il vettore verso il \textbf{centro del cubo adiacente}: rapporto $\sqrt{\frac{11}{3}}$.
\end{enumerate}

Queste sono le uniche possibilità compatibili con una simmetria cubica (le tre opzioni essendo cubico semplice, FCC e BCC). Poiché il reticolo reciproco è BCC, il \textbf{reticolo di Bravais diretto è FCC}.

\subsection{FCC come cubico semplice con base}



Per risolvere l'esercizio in modo sistematico, conviene descrivere il reticolo FCC come un \textbf{cubico semplice con base fittizia} di 4 atomi. Questo approccio presenta due vantaggi fondamentali:

\begin{enumerate}
    \item I \textbf{vettori fondamentali del reticolo reciproco} sono ortogonali (quelli del cubico semplice), mentre i vettori primitivi del BCC non lo sono.
    \item L'espressione del modulo di $\mathbf{K}$ è molto più semplice in termini degli indici di Miller del cubico semplice.
\end{enumerate}

\subsubsection{Vettori del reticolo reciproco}

Nel cubico semplice di lato $a$, i vettori del reticolo reciproco hanno la forma:
\begin{equation}
\mathbf{K} = \frac{2\pi}{a}\left(h\hat{\mathbf{i}} + k\hat{\mathbf{j}} + l\hat{\mathbf{k}}\right)
\end{equation}
con $h, k, l$ interi.

\subsubsection{Base FCC e regola di selezione}

La base fittizia dell'FCC nel cubico semplice è composta da 4 atomi equivalenti nelle posizioni:
\begin{equation}
\mathbf{d}_1 = (0, 0, 0), \quad \mathbf{d}_2 = \frac{a}{2}(0, 1, 1), \quad \mathbf{d}_3 = \frac{a}{2}(1, 0, 1), \quad \mathbf{d}_4 = \frac{a}{2}(1, 1, 0)
\end{equation}

Il fattore di struttura di questa base fittizia è non nullo solo quando $h$, $k$ e $l$ soddisfano la condizione:

\begin{equazione}{Regola di selezione FCC}
\boxed{S(\mathbf{K}) \neq 0 \iff h, k, l \text{ tutti pari o tutti dispari}}
\end{equazione}

Questa è la \textbf{regola di selezione} che trasforma il reticolo reciproco cubico semplice nel BCC, cioè nel reticolo reciproco dell'FCC.

\subsection{Etichettatura sistematica dei picchi}



La distanza interplanare nel Debye-Scherrer è:
\begin{equation}
d = \frac{\lambda}{2\sin\theta} = \frac{2\pi}{|\mathbf{K}|}
\end{equation}

Poiché $|\mathbf{K}| = \frac{2\pi}{a}\sqrt{h^2 + k^2 + l^2}$, si ottiene:

\begin{equation}
\boxed{d = \frac{a}{\sqrt{h^2 + k^2 + l^2}}}
\end{equation}

\begin{nota}
\textbf{Nota:} questa è la distanza reale tra i piani reticolari solo quando $h$, $k$ e $l$ sono coprimi. Se non sono coprimi, il fattore comune $n$ corrisponde all'ordine di molteplicità $n$ nella formula di Bragg ($2d\sin\theta = n\lambda$).
\end{nota}

E invertendo per il parametro reticolare:

\begin{equation}
\boxed{a = \frac{\lambda}{2\sin\theta}\sqrt{h^2 + k^2 + l^2}}
\end{equation}

\subsubsection{Lista degli indici di Miller}

Si elencano sistematicamente le terne $(hkl)$ con $h, k, l$ \textbf{tutti pari o tutti dispari}, ordinate per valore crescente di $h^2 + k^2 + l^2$. Le permutazioni sono equivalenti nel Debye-Scherrer (contano solo i moduli):

\begin{table}[h!]
\centering
\begin{tabular}{|c|c|c|}
\hline
Picco & $(hkl)$ & $h^2 + k^2 + l^2$ \\ \hline
1 & $(111)$ & $3$ \\
2 & $(200)$ & $4$ \\
3 & $(220)$ & $8$ \\
4 & $(311)$ & $11$ \\
5 & $(222)$ & $12$ \\
6 & $(400)$ & $16$ \\
7 & $(331)$ & $19$ \\
8 & $(420)$ & $20$ \\
9 & $(422)$ & $24$ \\
10 & $(333)/(511)$ & $27$ \\ \hline
\end{tabular}
\end{table}

\begin{nota}
\textbf{Perché $(100)$ non è ammesso?} Perché $1, 0, 0$ non sono tutti pari né tutti dispari. Il primo insieme ammesso con il minor valore di $h^2 + k^2 + l^2$ è $(111)$ con somma $3$, che precede $(200)$ con somma $4$.
\end{nota}



Il decimo picco presenta una particolarità notevole: \textbf{due famiglie di piani} diverse, $(333)$ e $(511)$, danno lo stesso valore $h^2 + k^2 + l^2 = 27$ e sono entrambe compatibili con la regola di selezione (tutti dispari). Entrambe producono lo \textbf{stesso picco} nel Debye-Scherrer, poiché il metodo delle polveri è sensibile solo al modulo di $\mathbf{K}$, non alla sua direzione.

\begin{nota}
\textbf{Nota sugli indici coprimi:} gli indici $(222)$ non sono coprimi; dividendo per il fattore comune $2$ si ottiene $(111)$. Ciò significa che il picco $(222)$ è la \textbf{riflessione di secondo ordine} ($n = 2$) dalla famiglia di piani $(111)$. Analogamente, $(333)$ è la riflessione di terzo ordine della stessa famiglia. Tuttavia, nel contesto del cubico semplice, li trattiamo come picchi distinti associati a vettori $\mathbf{K}$ diversi.
\end{nota}

\subsection{Calcolo del parametro reticolare $a$}



Per ciascun picco, il parametro reticolare si calcola come:
\begin{equation}
a = \frac{\lambda}{2\sin\theta}\sqrt{h^2 + k^2 + l^2}
\end{equation}
con $\lambda = 0{,}1542$ nm.



\begin{table}[h!]
\centering
\begin{tabular}{|c|c|c|c|c|c|}
\hline
Picco & $\theta$ [deg] & $\sin\theta$ & $(hkl)$ & $h^2+k^2+l^2$ & $a$ [nm] \\ \hline
1 & $18{,}475$ & $0{,}3169$ & $(111)$ & $3$ & $0{,}4214$ \\
2 & $21{,}455$ & $0{,}3658$ & $(200)$ & $4$ & $0{,}4215$ \\
3 & $31{,}150$ & $0{,}5173$ & $(220)$ & $8$ & $0{,}4216$ \\
4 & $37{,}320$ & $0{,}6063$ & $(311)$ & $11$ & $0{,}4217$ \\
5 & $39{,}320$ & $0{,}6337$ & $(222)$ & $12$ & $0{,}4215$ \\
6 & $47{,}030$ & $0{,}7317$ & $(400)$ & $16$ & $0{,}4215$ \\
7 & $52{,}875$ & $0{,}7973$ & $(331)$ & $19$ & $0{,}4215$ \\
8 & $54{,}890$ & $0{,}8180$ & $(420)$ & $20$ & $0{,}4215$ \\
9 & $63{,}645$ & $0{,}8961$ & $(422)$ & $24$ & $0{,}4215$ \\
10 & $71{,}885$ & $0{,}9504$ & $(333)/(511)$ & $27$ & $0{,}4215$ \\ \hline
\end{tabular}
\end{table}

I 10 valori ottenuti differiscono leggermente a causa delle incertezze numeriche. Il valore medio fornisce il \textbf{parametro reticolare della cella convenzionale FCC}:

\begin{equation}
\boxed{a \approx 0{,}4215 \text{ nm}}
\end{equation}

\subsection{Fattore di struttura con base reale}

Fino a questo punto si è trattato l'FCC come un cubico semplice con base \textbf{fittizia} di 4 atomi equivalenti. Il reticolo ha però formula chimica AB, dunque possiede una \textbf{base reale} di 2 atomi, in aggiunta alla base fittizia dell'FCC.

L'approccio più efficiente è il seguente:

\begin{enumerate}
    \item Si parte dai vettori $\mathbf{K}$ del reticolo reciproco del cubico semplice con la restrizione \textbf{$hkl$ tutti pari o tutti dispari} (che garantisce di avere i vettori del reticolo reciproco dell'FCC, cioè del BCC).
    \item Si calcola il fattore di struttura della \textbf{base reale} a 2 atomi con $\mathbf{d}_1 = (0, 0, 0)$ e $\mathbf{d}_2$ dato dall'esercizio.
\end{enumerate}

La base reale dell'FCC è costituita da 4 vettori di base. Ma poiché abbiamo già tenuto conto della base fittizia tramite la regola di selezione, il fattore di struttura da calcolare è solo quello della base reale a 2 atomi:
\begin{equation}
S_{\mathbf{K}} = \sum_{j=1}^{2} e^{-i\mathbf{K}\cdot\mathbf{d}_j} = 1 + e^{-i\mathbf{K}\cdot\mathbf{d}_2}
\end{equation}
dove il primo termine è sempre $1$ (poiché $\mathbf{d}_1 = \mathbf{0}$).

\begin{nota}
\textbf{Punto chiave:} se i due atomi sono \textbf{diversi} (A $\neq$ B), il fattore di struttura include i fattori di forma atomico $f_A$ e $f_B$, che sono in generale diversi. In tal caso è \textbf{molto improbabile} che $S_{\mathbf{K}}$ si annulli esattamente: tutti i picchi saranno visibili, ma con intensità \textbf{modulata}. La soppressione completa avviene \textbf{solo} quando gli atomi sono identici (A = B).
\end{nota}

\subsubsection{Caso 1: $\mathbf{d}_2 = a\left(\frac{1}{4}, \frac{1}{4}, \frac{1}{4}\right)$}

Si calcola il prodotto scalare:
\begin{equation}
\mathbf{K} \cdot \mathbf{d}_2 = \frac{2\pi}{a}(h, k, l) \cdot a\left(\frac{1}{4}, \frac{1}{4}, \frac{1}{4}\right) = \frac{\pi}{2}(h + k + l)
\end{equation}

Il fattore di struttura diventa:
\begin{equation}
S_{\mathbf{K}} = 1 + e^{-i\frac{\pi}{2}(h+k+l)}
\end{equation}

Questo si annulla quando $e^{-i\frac{\pi}{2}(h+k+l)} = -1$, ossia quando $\frac{h+k+l}{2}$ è un \textbf{numero intero dispari}.

Applicando questa regola a ciascun picco:

\begin{table}[h!]
\centering
\begin{tabular}{|c|c|c|c|}
\hline
$(hkl)$ & $h+k+l$ & $\frac{h+k+l}{2}$ & Scompare? \\ \hline
$(111)$ & $3$ & $1{,}5$ (non intero) & No \\
$(200)$ & $2$ & $1$ (dispari) & \textbf{Sì} \\
$(220)$ & $4$ & $2$ (pari) & No \\
$(311)$ & $5$ & $2{,}5$ (non intero) & No \\
$(222)$ & $6$ & $3$ (dispari) & \textbf{Sì} \\
$(400)$ & $4$ & $2$ (pari) & No \\
$(331)$ & $7$ & $3{,}5$ (non intero) & No \\
$(420)$ & $6$ & $3$ (dispari) & \textbf{Sì} \\
$(422)$ & $8$ & $4$ (pari) & No \\
$(333)/(511)$ & $9/7$ & $4{,}5/3{,}5$ (non interi) & No \\ \hline
\end{tabular}
\end{table}

\begin{equation}
\boxed{\text{Picchi soppressi: } (200),\; (222),\; (420)}
\end{equation}

\subsubsection{Caso 2: $\mathbf{d}_2 = a\left(\frac{1}{2}, \frac{1}{2}, \frac{1}{2}\right)$}

Il calcolo è analogo:
\begin{equation}
\mathbf{K} \cdot \mathbf{d}_2 = \frac{2\pi}{a}(h, k, l) \cdot a\left(\frac{1}{2}, \frac{1}{2}, \frac{1}{2}\right) = \pi(h + k + l)
\end{equation}

Il fattore di struttura è:
\begin{equation}
S_{\mathbf{K}} = 1 + e^{-i\pi(h+k+l)}
\end{equation}

Questo si annulla quando $e^{-i\pi(h+k+l)} = -1$, ossia quando $h + k + l$ è \textbf{dispari}.

\begin{table}[h!]
\centering
\begin{tabular}{|c|c|c|c|}
\hline
$(hkl)$ & $h+k+l$ & Dispari? & Scompare? \\ \hline
$(111)$ & $3$ & Sì & \textbf{Sì} \\
$(200)$ & $2$ & No & No \\
$(220)$ & $4$ & No & No \\
$(311)$ & $5$ & Sì & \textbf{Sì} \\
$(222)$ & $6$ & No & No \\
$(400)$ & $4$ & No & No \\
$(331)$ & $7$ & Sì & \textbf{Sì} \\
$(420)$ & $6$ & No & No \\
$(422)$ & $8$ & No & No \\
$(333)/(511)$ & $9/7$ & Sì & \textbf{Sì} \\ \hline
\end{tabular}
\end{table}

\begin{equation}
\boxed{\text{Picchi soppressi: } (111),\; (311),\; (331),\; (333)/(511)}
\end{equation}


Le due scelte della base producono \textbf{pattern di diffrazione completamente diversi}:

\begin{table}[h!]
\centering
\begin{tabular}{|c|c|c|}
\hline
$(hkl)$ & $\mathbf{d}_2 = a(\frac{1}{4},\frac{1}{4},\frac{1}{4})$ & $\mathbf{d}_2 = a(\frac{1}{2},\frac{1}{2},\frac{1}{2})$ \\ \hline
$(111)$ & Visibile & \textbf{Soppresso} \\
$(200)$ & \textbf{Soppresso} & Visibile \\
$(220)$ & Visibile & Visibile \\
$(311)$ & Visibile & \textbf{Soppresso} \\
$(222)$ & \textbf{Soppresso} & Visibile \\
$(400)$ & Visibile & Visibile \\
$(331)$ & Visibile & \textbf{Soppresso} \\
$(420)$ & \textbf{Soppresso} & Visibile \\
$(422)$ & Visibile & Visibile \\
$(333)/(511)$ & Visibile & \textbf{Soppresso} \\ \hline
\end{tabular}
\end{table}

Cambiando la base, il pattern di diffrazione cambia radicalmente. Dunque, anche un esperimento di Debye-Scherrer (che perde l'informazione sulla direzione di $\mathbf{K}$) è sufficiente per \textbf{distinguere} tra le due strutture cristalline.

Quando gli atomi \textbf{non sono identici} (base reale con $A \neq B$), i picchi non scompaiono completamente ma risultano solo \textbf{attenuati} rispetto ai picchi vicini, a causa della soppressione parziale dovuta al fattore di struttura con fattori di forma atomico diversi.

\section{Il reticolo reciproco e la base}

Questo è un punto concettuale fondamentale, sottolineato ripetutamente dal professore. Le posizioni atomiche nel cristallo sono della forma:
\begin{equation}
\mathbf{r}_c = \mathbf{R} + \mathbf{d}_j
\end{equation}

\begin{itemize}
    \item $\mathbf{R}$ è un vettore del reticolo di Bravais $\to$ \textbf{insieme infinito e periodico} $\to$ definisce il reticolo reciproco.
    \item $\{\mathbf{d}_j\}$ ($j = 1, \ldots, p$) è la base $\to$ \textbf{insieme finito e non periodico} $\to$ \textbf{non ammette} un reticolo reciproco.
\end{itemize}


\textbf{Punto fondamentale:} il reticolo reciproco è definito \textbf{esclusivamente} dal reticolo di Bravais $\mathbf{R}$, non dalla base. La base, essendo un insieme finito di vettori, non possiede periodicità e quindi non ha senso definirne un reticolo reciproco. La base interviene solo nella \textbf{modulazione dell'intensità} dei picchi tramite il fattore di struttura $S(\mathbf{K})$.


Conseguenze pratiche:
\begin{itemize}
    \item Se gli atomi della base sono \textbf{diversi} (base reale), \textbf{tutti} i picchi del reticolo reciproco dell'FCC sono visibili, indipendentemente dalla base. La base modula solo l'intensità.
    \item Se gli atomi sono \textbf{identici}, il fattore di struttura può annullarsi per certi $\mathbf{K}$, sopprimendo completamente alcuni picchi.
    \item Per determinare il tipo di reticolo di Bravais, la base è \textbf{irrilevante}: basta analizzare le posizioni dei picchi.
\end{itemize}

\newpage
\section{Introduzione alla dinamica reticolare}


Un cristallo ideale è un \textbf{arrangiamento periodico} di atomi o molecole. Molto era noto sulla regolarità dei cristalli già prima della scoperta della struttura atomica della materia: i mineralogisti avevano osservato che gli angoli tra le facce dei cristalli erano sempre gli stessi per un dato minerale.

Sappiamo oggi che la materia, composta di atomi e molecole, a temperature sufficientemente basse e densità sufficientemente alte forma una \textbf{fase solida}. Se questa è una fase di equilibrio termodinamico, il solido è un \textbf{cristallo}. Esistono anche fasi solide metastabili, i \textbf{solidi amorfi}, ma la fase di equilibrio è cristallina.

Nel modello ideale a $T = 0$ (stato fondamentale), gli atomi occupano posizioni di equilibrio perfettamente periodiche. Non appena si aumenta la temperatura, gli atomi iniziano a \textbf{oscillare} attorno alle loro posizioni di equilibrio.



Dal punto di vista fondamentale, il sistema è composto da \textbf{ioni} (nuclei con elettroni di core) e \textbf{elettroni} di valenza. In linea di principio, si dovrebbe risolvere l'equazione di Schrödinger per tutte le particelle:
\begin{equation}
H\psi(\{\mathbf{r}_e\}, \{\mathbf{R}_N\}) = E\,\psi(\{\mathbf{r}_e\}, \{\mathbf{R}_N\})
\end{equation}
dove $\{\mathbf{r}_e\}$ sono le coordinate elettroniche e $\{\mathbf{R}_N\}$ quelle nucleari. Questo è un problema \textbf{enormemente complesso}: l'unico sistema risolvibile esattamente è la molecola $\text{H}_2^+$ (due ioni e un solo elettrone). Un cristallo può essere pensato come una \textbf{molecola enorme}, e molto del linguaggio utilizzato per i solidi proviene dalla fisica molecolare.

\subsection{L'approssimazione di Born-Oppenheimer}

Per affrontare il problema, si ricorre all'\textbf{approssimazione di Born-Oppenheimer}, nota dalla fisica molecolare. L'idea fondamentale è che gli \textbf{elettroni sono molto più leggeri degli ioni} ($m \ll M$), e quindi si muovono molto più velocemente.

Nell'Hamiltoniana, i termini cinetici sono:
\begin{itemize}
    \item Per gli ioni: $\frac{p^2}{2M}$ (dove $M$ è la massa ionica)
    \item Per gli elettroni: $\frac{p^2}{2m}$ (dove $m$ è la massa elettronica, $m \ll M$)
\end{itemize}

L'approssimazione consiste nel \textbf{separare le due scale di moto}:

\begin{enumerate}
    \item Si risolve il \textbf{problema elettronico} assumendo che le posizioni dei nuclei $\{\mathbf{R}_N\}$ siano \textbf{fissate} (parametri, non variabili dinamiche). Si ottengono gli autovalori elettronici $E_n(\{\mathbf{R}_N\})$, che dipendono parametricamente dalle posizioni nucleari.
    \item Gli autovalori elettronici dello stato fondamentale $E_0(\{\mathbf{R}_N\})$ fungono da \textbf{potenziale efficace} per il moto dei nuclei: i nuclei si muovono in un potenziale generato dagli elettroni.
\end{enumerate}

Questo potenziale efficace, come funzione delle posizioni nucleari, presenta dei \textbf{minimi periodici} (in un solido). Questi minimi corrispondono alle \textbf{posizioni di equilibrio} degli atomi nel cristallo.

\subsection{I fononi: quanti del suono}

Se si espande il potenziale efficace attorno ai minimi (approssimazione armonica, valida per piccole oscillazioni), si ottiene un sistema di \textbf{oscillatori armonici accoppiati}. La quantizzazione di questo sistema produce livelli energetici discreti:
\begin{equation}
E_n = \left(n + \frac{1}{2}\right)\hbar\omega
\end{equation}

I \textbf{quanti} di queste oscillazioni sono chiamati \textbf{fononi}, in analogia con i \textbf{fotoni} (quanti della luce). Le proprietà fondamentali dei fononi sono:

\begin{itemize}
    \item Sono \textbf{bosoni}, obbediscono alla statistica di \textbf{Bose-Einstein}.
    \item Sono i quanti di un \textbf{campo} (il campo di deformazione del solido): le coordinate del campo sono le deviazioni degli atomi dalle posizioni di equilibrio.
    \item A \textbf{basse energie}, le oscillazioni sono genuinamente quantistiche.
    \item Ad \textbf{alte energie} (molti quanti), le oscillazioni si avvicinano alle \textbf{onde sonore classiche}, che possiamo percepire.
\end{itemize}

